\section{Presentación del invitado}

Jiayi Weng (翁家翌) ingresó en 2016 a la licenciatura en Ciencias de la Computación de la Universidad de Tsinghua, en 2020 partió a la Universidad Carnegie Mellon (CMU) a cursar la maestría y en 2022 se incorporó a OpenAI. Es uno de los constructores centrales de \textbf{Post-Training RL Infrastructure} dentro de OpenAI: detrás de cada modelo grande que OpenAI ha publicado, desde ChatGPT (GPT-3.5) y GPT-4o hasta GPT-5, está su nombre.

Antes de unirse a OpenAI, Jiayi Weng ya había generado un impacto amplio a través de proyectos de código abierto: liberó en Tsinghua todas sus tareas y materiales de curso para romper la asimetría de información; creó el framework de aprendizaje por refuerzo \textbf{Tianshou (天授)}, que se convirtió en uno de los frameworks ligeros más populares de la comunidad de RL; y durante la pandemia desarrolló el sistema gratuito de consulta de visas \textbf{tuixue.online}, que acumuló millones de visitas.

Esta edición del podcast WhynotTV es una conversación a fondo de más de dos horas que arranca con la infancia de Jiayi Weng y abarca sus estudios, el código abierto, sus decisiones profesionales, la forma de trabajar dentro de OpenAI, y sus reflexiones filosóficas sobre el futuro de la IA y el sentido de la vida.

\begin{knowledgebox}{Contexto de la entrevista}
El outline de esta edición del podcast lo preparó el conductor con la función Deep Research de GPT-5, y Jiayi Weng es precisamente uno de los desarrolladores centrales detrás de GPT-5. El conductor lo llamó «un ciclo maravilloso».
\end{knowledgebox}

\subsection{Resumen de la sección}

Las tres palabras clave de Jiayi Weng son: \textbf{aprendizaje por refuerzo, Post-Training, Infra}. Pero su historia va mucho más allá de lo técnico: es alguien que hace «caridad» con código, que cree en el código abierto y en la igualdad de acceso a la información, y que mantiene un pensamiento independiente en el centro de la tormenta mundial de la IA.

